% ECONOMICS_PROBLEM: {"id": "D63-70", "title": "Polynomial-time construction with subquadratic envy-free subsidies", "jel_code": "D63", "source_id": "OP-0593"}
% FIRST_POSTED: 2026-10-09
% ATTEMPT_STATUS: unresolved
% ENTRY_KIND: research_attempt
% !TeX program = pdflatex
\documentclass[11pt,a4paper]{article}
\usepackage[T1]{fontenc}
\usepackage[utf8]{inputenc}
\usepackage{lmodern}
\usepackage[margin=24mm]{geometry}
\usepackage{amsmath,amssymb,amsthm,mathtools}
\usepackage{booktabs,longtable,array,enumitem,xurl,hyperref,listings,microtype}
\hypersetup{colorlinks=true,linkcolor=blue,urlcolor=blue,citecolor=blue,pdftitle={Research proof audit}}
\setlength{\parindent}{0pt}
\setlength{\parskip}{5pt}
\setlength{\emergencystretch}{3em}
\setlist{nosep,leftmargin=2em}
\newtheorem{theorem}{Theorem}[section]
\newtheorem{lemma}[theorem]{Lemma}
\newtheorem{proposition}[theorem]{Proposition}
\newtheorem{corollary}[theorem]{Corollary}
\theoremstyle{remark}\newtheorem{remark}[theorem]{Remark}
\newcommand{\R}{\mathbb R}
\newcommand{\N}{\mathbb N}
\newcommand{\E}{\mathbb E}
\newcommand{\Var}{\operatorname{Var}}
\newcommand{\ind}{\mathbf 1}
\newcommand{\EF}{\mathrm{EF}}
\newcommand{\EFM}{\mathrm{EFM}}
\newcommand{\PO}{\mathrm{PO}}
\newcommand{\MMS}{\mathrm{MMS}}
\DeclareMathOperator*{\argmax}{arg\,max}
\DeclareMathOperator*{\argmin}{arg\,min}
\lstset{basicstyle=\ttfamily\footnotesize,breaklines=true,columns=fullflexible,keepspaces=true,frame=single}
\begin{document}
\begin{center}
{\LARGE\bfseries Polynomial computation of subquadratic subsidies}\par\medskip
{\large D63-70.tex}\par
Research attempt and adversarial audit --- 9 October 2026
\end{center}
\label{problem:OP-0593}
\label{audited:ECON-006}
\label{audited:conventions}
\textbf{Tools.} Web browsing: YES. Computation: YES (Python, exact integer/rational checks, and explicitly identified numerical diagnostics).
\textbf{Status convention.} A full proof of a restricted theorem does not resolve the full input. Literature results are marked as imported; no novelty or exhaustive current-literature certification is claimed. Computational searches certify only their stated finite domains.

\section{FORMAL RESTATEMENT}
For every $n\ge1$, finite $M$ with $m=|M|$, and family of valuations $v_i:2^M\to\mathbb Q$ satisfying
\[
v_i(\varnothing)=0,\qquad 0\le v_i(S\cup\{g\})-v_i(S)\le1\quad(g\notin S),
\]
a value oracle returns an exact rational of encoding length at most the supplied bound $L$. Does there exist an absolute $C$ and an algorithm using $\operatorname{poly}(n,m,L)$ oracle calls and bit operations that outputs a complete allocation and nonnegative rational payments with
\[
v_i(A_i)+p_i\ge v_i(A_j)+p_j\quad(\forall i,j),\qquad
\sum_i p_i\le C n^{3/2}\sqrt{\log(4n)}?
\]
Logarithms are natural. The same constant and polynomial must work for all inputs. A randomized variant must specify probability of success and the runtime convention; our restricted algorithms are deterministic. Approximate envy, demand-oracle access, explicit exponential tables and succinct circuits are different models. Empty bundles and zero valuations are allowed. A feasible output must have polynomial encoding length; exact irrational output is not authorized.

\section{QUICK LITERATURE/CONTEXT CHECK}
The previously inaccessible source is now readable. Its Theorem 1 gives an existential bound $5\sqrt2(n-1)^{3/2}\sqrt{\log(4n)}$ for nonnegative bundle values under the absolute marginal bound; monotone goods satisfy those hypotheses. Section 6 asks about efficient computation \cite{subq}. Thus existence at the requested scale is not the missing step here.

\section{ATTACK PLAN}
The proof track isolates polynomial-time payment computation and searches for partitions in tractable subdomains. The construction track tests positive-cycle obstructions and small nonadditive examples. A full enumeration supplies an exact benchmark, but is not confused with a polynomial value-oracle algorithm.

\begin{center}\begin{tabular}{p{.25\linewidth}p{.66\linewidth}}\toprule Tool&Role\\\midrule
Difference constraints&Characterize exactly which partitions admit subsidies.\\
Longest-path potentials&Give minimal nonnegative payments and a certificate.\\
Finite welfare-maximizing assignment&Ensures feasibility for at least one bundle assignment.\\
Prefix crossing&Produces the sharp two-agent bound.\\
Greedy load balancing&Solves identical monotone valuations.\\
Oracle/bit-complexity accounting&Separates evaluation of a certificate from finding it.\\
Small Lipschitz set-function search&Tests nonadditive and mixed-sign boundary cases.\\
Topological/discrepancy methods&Possible route to global bounds, but not a polynomial algorithm by themselves.\\\bottomrule\end{tabular}\end{center}

\section{WORK}

\subsection{A complete payment certificate for a fixed partition}
\begin{lemma}[Difference constraints, exact feasibility and attainment]\label{lem:graph}
For a complete allocation $A$, put $w_{ij}=v_i(A_j)-v_i(A_i)$. There exist nonnegative subsidies making $A$ envy-free if and only if every directed cycle of this complete weighted graph has nonpositive total weight. In that case the coordinatewise smallest nonnegative feasible vector is
\[
p_i^*=\max\bigl\{0,\ \text{total weights of directed paths starting at }i\bigr\},
\]
where the maximum may be restricted to simple paths. In particular $\min_i p_i^*=0$. For every finite valuation instance there is at least one envy-freeable complete allocation, and the infimum defining $s(v)$ is attained.
\end{lemma}
\begin{proof}
The envy inequalities are $p_i\ge w_{ij}+p_j$. Summing around a cycle gives zero on the payment side, so positive cycles are impossible. Conversely, when no cycle is positive, deleting a repeated-vertex segment from a walk does not decrease its weight. The largest walk weight is therefore attained by a simple path of length at most $n-1$, including the empty path of weight zero. Prepending edge $i\to j$ to a maximizing path from $j$ and deleting any resulting nonpositive cycles proves $p_i^*\ge w_{ij}+p_j^*$. Thus $p^*$ is feasible and nonnegative.

Any nonnegative feasible $p$ dominates every path weight starting at $i$: sum the edge inequalities along that path and use nonnegativity of its final payment. Hence $p_i\ge p_i^*$ for every $i$. If the minimum coordinate of $p^*$ were positive, subtracting it from all coordinates would give a strictly smaller nonnegative feasible vector, a contradiction.

To obtain some envy-freeable allocation, begin with any $n$ bundle occurrences, including empty occurrences, forming a partition. Choose an assignment of those bundles to the agents that maximizes total value over the finite set of $n!$ assignments. A positive cycle in its envy graph would reassign the bundles cyclically and strictly increase that total, an impossibility. Finally there are finitely many complete allocations. Each feasible one has the attained minimum total $\sum_i p_i^*$, and at least one is feasible; taking the smallest of finitely many such totals proves attainment of $s(v)$.
\end{proof}

\begin{corollary}[An item-dependent bound, not a uniform theorem]\label{cor:mdep}
If $v_i(\varnothing)=0$ and every absolute one-item marginal is at most one, then some envy-free allocation has total subsidy at most $(n-1)^2m$, where $m=|M|$.
\end{corollary}
\begin{proof}
Adding the items of a bundle one at a time gives $|v_i(S)|\le |S|$. Distinct allocated bundles are disjoint, so $|w_{ij}|\le|A_i|+|A_j|\le m$. On an envy-freeable allocation, every simple path has at most $n-1$ edges, whence $p_i^*\le(n-1)m$. At least one coordinate is zero, giving the bound. It is zero when $n=1$ or $m=0$.
\end{proof}
The supremum over all finite $M$ makes this estimate insufficient for any proposed bound depending only on $n$.

\begin{lemma}[Exact evaluation with polynomial-size payments]\label{lem:compute}
For a specified allocation with rational oracle values of bit length at most $L$, envy-freeability and the minimum payments can be computed with $n^2$ value queries and $O(n^3)$ rational arithmetic operations. The output bit length is polynomial in $n,m,L$.
\end{lemma}
\begin{proof}
Set $p^{(0)}=0$ and iterate
\[
p_i^{(k+1)}=\max\{p_i^{(k)},\max_j(w_{ij}+p_j^{(k)})\}.
\]
An induction on $k$ identifies $p_i^{(k)}$ as the greatest weight of a walk starting at $i$ of length at most $k$, also allowing zero. If $p^{(n)}=p^{(n-1)}$, this vector satisfies all constraints and equals the minimum vector. If an entry increases, some length-$n$ walk improves on all shorter walks; it must contain a positive cycle, since deleting only nonpositive cycles could not lower its weight. Thus the allocation is infeasible. The recurrence has $n$ rounds of $n^2$ comparisons/additions. A path sums at most $n$ differences of input rationals. Forming common denominators by multiplication bounds its encoding length by $O(nL+\log n)$; each maximum selects one such rational. Query arguments are $m$-bit masks. These facts give polynomial bit complexity, not just a real-arithmetic claim.
\end{proof}

\subsection{A sharp theorem for two agents, even without monotonicity}
\begin{theorem}\label{thm:two}
For $n=2$, any normalized valuations with absolute one-item marginals at most one have an envy-free allocation with total subsidy at most one. This can be found with $O(m)$ value queries. The worst-case constant one is sharp.
\end{theorem}
\begin{proof}
Order the items arbitrarily and let $S_k$ be the first $k$, $0\le k\le m$. Write $d_i(S)=v_i(S)-v_i(M\setminus S)$. Consecutive $d_1(S_k)$ differ in absolute value by at most two, because an item is added on one side and removed on the other. The first and last are $-v_1(M)$ and $v_1(M)$. Hence some $k$ has $|d_1(S_k)|\le1$: otherwise the sequence must cross from a number below $-1$ to a number above $1$, or in the reverse direction, in one step of absolute size greater than two. This also covers $v_1(M)=0$ and $m=0$ by choosing $k=0$.

Fix this bipartition $S,T$. Write $d_i=v_i(S)-v_i(T)$. If $d_1\ge d_2$, assign $S$ to agent 1 and $T$ to agent 2. The constraints on the payment difference are
\[
d_2\le p_2-p_1\le d_1.
\]
Choose the element of this closed interval nearest zero; its absolute value is the minimum possible total nonnegative subsidy, by placing that amount on the appropriate agent and zero on the other. If $d_1<d_2$, reverse the bundles and apply the same argument to $-d_1,-d_2$. The resulting minimum total is zero if $d_1,d_2$ straddle zero, and otherwise $\min\{|d_1|,|d_2|\}$. It is therefore at most one. Scanning the prefixes requires $O(m)$ queries and the final pair of values for agent 2; all arithmetic is exact for rational inputs.

For sharpness, use one good valued at one by both agents. Whichever receives it, the other needs a subsidy exactly one larger than the recipient's. Nonnegative payments therefore have total at least one, achieved by paying the nonrecipient one and the recipient zero.
\end{proof}
For monotone goods the sequence $d_1(S_k)$ is nondecreasing, so the sign crossing can instead be located by binary search with $O(\log(m+1))$ queries. This improvement is not asserted for arbitrary nonmonotone functions.

\subsection{Identical monotone valuations: a linear constructive bound}
\begin{theorem}\label{thm:identical}
If all agents have the same normalized monotone valuation $v$ with marginals in $[0,1]$, there is an envy-free allocation with $p_i\in[0,1]$ and total subsidy at most $n-1$, obtainable with polynomially many value queries and bit operations for rational data.
\end{theorem}
\begin{proof}
Initially all bundles are empty and their values are zero. Process the items in any order; give the next item to a bundle of minimum current value. If the old minimum is $a$ and all bundle values are in $[a,a+1]$, the receiving bundle's new value is also in $[a,a+1]$, by the marginal bounds, and every other value is unchanged. Thus by induction the spread of the final values $\ell_i=v(A_i)$ is at most one. Let $H=\max_i\ell_i$ and $p_i=H-\ell_i$. Every compared bundle-payment pair has the same value $H$ to every agent. These payments are nonnegative, at most one, and at least one is zero, so their sum is at most $n-1$. Only the changed bundle needs a new oracle query at each step. Minimum selection and rational subtraction have polynomial bit complexity. Equality of the agents' valuations is a promise for this subcase, not something proved by finitely many queries.
\end{proof}

\begin{proposition}[A universal linear lower example]\label{prop:lower}
For every $n\ge1$, the instance consisting of one good valued at one by every agent has $s(v)=n-1$.
\end{proposition}
\begin{proof}
If agent $r$ receives the good, envy-freeness between $r$ and any other agent $i$ imposes both $p_i\ge1+p_r$ and $p_i\le1+p_r$. Thus $p_i=1+p_r$ for all $i\ne r$. The minimum nonnegative vector has $p_r=0$ and all other coordinates one, giving total $n-1$. This includes $n=1$.
\end{proof}

\subsection{A complete but exponential algorithm}
\begin{proposition}\label{prop:exhaustive}
For every rational input, exhaustive search finds an exact minimum-total-subsidy allocation with rational payments. Using the imported existence theorem, this output meets the requested numerical bound with $C=5\sqrt2$, but the algorithm is not polynomial in $m$.
\end{proposition}
\begin{proof}
Query every $v_i(S)$, at most $n2^m$ values, and enumerate all $n^m$ assignments of items to agents. For each, apply Lemma~\ref{lem:compute}; discard infeasible partitions and retain the smallest exact total over feasible ones. Lemma~\ref{lem:graph} proves feasibility and attainment, so the procedure returns a true optimum. The imported existence theorem applies because normalization and monotonicity imply $v_i(S)\ge0$ and the required absolute marginal bound holds. The optimum is no larger than its bound, hence no larger than $5\sqrt2 n^{3/2}\sqrt{\log(4n)}$. No irrational logarithm needs to be computed to run this minimization: all objective comparisons use exact rational sums. The explicitly exponential query count and enumeration preclude the claimed polynomial-time conclusion.
\end{proof}

\section{VERIFICATION}

The executable checks generated normalized integer set functions of the form
\[
v(S)=\min_r\{c_r+|S\triangle K_r|\}-\min_r\{c_r+|K_r|\}.
\]
Each cone is 1-Lipschitz for the Hamming metric, and the pointwise minimum is also 1-Lipschitz: using a minimizing index at one endpoint gives one direction of the inequality, and interchanging endpoints gives the other. The program additionally checked every one-item marginal explicitly. With seed 20261009, it tested 200 two-agent/five-item instances and 200 three-agent/four-item instances. It enumerated every complete allocation, computed exact minimum payments by the recurrence, and found largest observed optima of one in both batches. These are nonuniform small samples, not worst-case asymptotic bounds; many generated valuations are not monotone.
\begin{lstlisting}
for each sampled valuation table v:
    best = infinity
    for each complete allocation A:
        w[i,j] = v[i,A[j]] - v[i,A[i]]
        p = longest-walk recurrence from zero
        if p passes all envy constraints: best = min(best,sum(p))
    retain exact best and a witnessing allocation
\end{lstlisting}
The one-good lower example was verified algebraically for every $n$, not inferred from this sample. Boundary checks include $n=1$, no items, zero values, zero-weight cycles, and exact rational payment differences. No positive-cycle partition is declared feasible merely because one or more agents could be paid a large amount: a positive cycle makes every payment vector impossible.

The arbitrary-table tests stress the payment evaluator on a larger domain; they are not themselves evidence that the original monotone oracle problem is hard. The identical-valuation algorithm uses only the promised common oracle values. The longest-path computation proves that \emph{once an appropriate partition is known}, irrational payments and unbounded output length are not barriers. It does not bound the search for that partition. A demand-oracle subroutine is nowhere substituted for a value query.

\section{FINAL}
\textbf{LABEL: UNRESOLVED.}
The target polynomial value-oracle algorithm is neither proved nor ruled out. Exact polynomial certificate evaluation, the sharp two-agent case, an all-$n$ identical-valuation algorithm and a correct exponential benchmark are fully established.

\textbf{Exact first gap.} No polynomial-query method is provided for finding, among all complete partitions, one whose minimum payment vector has total $O(n^{3/2}\sqrt{\log n})$ under heterogeneous monotone valuations.

\textbf{Three next moves.} (1) Find a polynomially evaluable separation or descent oracle for a partition potential tied to minimum subsidies. (2) Turn a fractional or topological existence certificate into an explicit finite rational search region with polynomial complexity. (3) Build a value-oracle indistinguishability family whose low-subsidy partitions reveal a hidden bundle, and prove a genuine query lower bound rather than treating nonconstructivity as hardness.

\textbf{Minimal hard instance.} At least three nonidentical valuations are needed to escape the two elementary algorithms here; sharply context-dependent marginal values are a plausible obstruction. No lower-bound family is certified by the small random search.

\section{Completion Estimate}
20\%

Exact payment computation and important restricted cases are complete; the main heterogeneous value-oracle search problem remains unresolved. The existing existential result is credited separately.

This is a subjective estimate of the fraction of the \emph{whole requested mathematical scope} addressed by this report, including the named variants. It is not a probability of correctness, an objectively measured fraction of a proof, or an estimate that the remaining work takes the complementary fraction of the time. Only the eleven supplied files are assessed; no missing rank-1--200 statements are inferred.

\begin{thebibliography}{99}
\bibitem{subq} M. Dupr\'e la Tour and M. Suzuki. \emph{Subquadratic Subsidies for Nonnegative or Nonpositive Valuations}. arXiv:2609.08272v1, 8 September 2026, Theorem 1 and Section 6. \url{https://arxiv.org/html/2609.08272v1}.
\end{thebibliography}
\end{document}
